\section{Introduction}
\label{sec:intro}

Much of experimental biology is a search for interventions that move a cell from one state
to another: differentiating a stem cell, repolarizing a T cell, reverting a diseased
phenotype, or---in drug discovery---knocking down a target to reverse a pathogenic program.
Before committing reagents to such a program, the decisive question is not \emph{what will
this perturbation do} but \emph{can the desired change be achieved at all with the
perturbations available, and if not, what is missing}. That question is rarely asked
computationally, and its economic weight is large: only about one in ten clinical
programs succeeds and many late-stage failures involve insufficient efficacy. That clinical
context motivates better early evidence, but a transcriptomic screen cannot by itself explain
or prevent clinical attrition. Human genetic support and directly measured target effects do,
however, both argue for grounding target choice in real perturbation data rather than inferred
models \citep{Minikel2024,Nelson2015,Harrison2016,%
BegleyEllis2012,Prinz2011}.

The computational tools available to answer it fall into two camps that, as we document in
Section~\ref{sec:related}, do not overlap. Perturbation-response models---from
gene-regulatory-network simulators to single-cell foundation models---are increasingly
trained on real screens but are forward predictors: given a perturbation they return a
predicted response, and a predictor structurally always returns \emph{some} answer, so it
cannot certify that a goal is out of reach. Network-control and Boolean-attractor methods
do reason about reachability and can carry controllability guarantees, but those guarantees
are properties of a hand-built or inferred model, only as trustworthy as the network
assumed. Neither camp answers the feasibility question from measured effects. This gap
matters all the more because a wave of 2025--2026 benchmarks finds that current
perturbation predictors, foundation models included, frequently fail to beat mean or linear
baselines on genuinely unseen perturbations \citep{AhlmannEltze2025,Wong2025,Kernfeld2025,%
Atti2025}: the field's own evidence argues for an output weaker and more defensible than
predicting an unseen response.

\paragraph{The convex-cone move.}
We take the opposite starting point from model building: we do not infer a network at all.
A genome-scale CRISPRi Perturb-seq screen \citep{Zhu2025} directly measures, for each
perturbation--condition profile, the effect vector it produces in expression
space. We treat these measured vectors as a fixed dictionary and ask a geometric question:
does the desired cell-state shift lie in the set of shifts the dictionary can produce?
Two modelling choices make this a convex-geometry problem. A measured knockdown effect may be
applied with non-negative weight but not reversed into an unmeasured ``negative knockdown'';
and several knockdowns are approximated, to first order, by the sum of their single effects.
Individual effect vectors may raise some transcripts and lower others. The set of directions
allowed by this model is therefore the \emph{non-negative conic hull} of
the measured effect vectors, and asking whether a target is reachable is asking whether it
lies inside that model-relative cone---a question non-negative least squares (NNLS) answers exactly
\citep{LawsonHanson1974}. This reframing turns an open-ended design search into a decision
with a checkable solution on both sides. When the target is inside the cone, the NNLS weights
rank a feasible mixture and a greedy spectrum supplies compact candidate panels; neither is a
global minimum-cardinality proof. When it is outside, convex duality supplies a separating
hyperplane---a Farkas/KKT certificate \citep{Farkas1902,KuhnTucker1951}---whose full vector
certifies the mismatch. Positive coordinates rank readouts under-delivered at the closest cone
point. They motivate a CRISPRa or de-repression arm, but do not prove that activating any
individual readout will cause the desired state.

\paragraph{Scope.}
The oracle is a decision instrument, not an oracle of biological truth: it certifies
feasibility \emph{relative to the measured dictionary} under an explicit additivity model, and
its verdicts are hypotheses for the bench, not validated targets. The regulators it surfaces in
T cells and in the CEBPA program are already reported by the source screens
\citep{Zhu2025,Norman2019} and enter here only as positive controls; the novelty we claim is
the measured-effect, certificate-carrying reachability \emph{decision} itself, not the biology
it recovers.

\paragraph{Contributions.}
This paper makes four contributions.
\begin{enumerate}
\item \textbf{A feasibility oracle for cell-state engineering.} We formalize reachability of
a target cell-state direction as membership in the convex cone of measured perturbation effect
vectors, and solve it with NNLS so that a verdict, a ranked mixture, a greedy sparse panel,
and---on infeasibility---a Farkas/KKT dual certificate follow from one convex program
(Section~\ref{sec:methods}).
\item \textbf{A validated headline result.} On primary human
CD4\textsuperscript{+} T cells the oracle finds T\textsubscript{H}2$\rightarrow$%
T\textsubscript{H}1 partly reachable by knockdown---far above a shuffled-target null but well
below full reachability---with a signed loss-/gain-of-function decomposition and canonical
regulators recovered as positive controls (Section~\ref{sec:results-headline}).
\item \textbf{Generality across transitions and datasets.} A 12-cell atlas shows knockdown
is never the majority modality across four transitions and three activation states, and the
identical operator runs without retuning on a K562 CRISPRa combinatorial screen, where
we also quantify the additivity approximation directly on measured doubles
(Sections~\ref{sec:results-atlas}--\ref{sec:results-transfer}).
\item \textbf{An exact positioning against the literature.} Over 92 methods spanning
2011--2026, we show that measured-effect grounding and feasibility-with-certificate are held
by disjoint method sets whose intersection is empty, and translate the resulting verdicts
into a GO/REDIRECT/STOP triage over 102 real target nominations
(Sections~\ref{sec:results-landscape}--\ref{sec:results-impact}).
\end{enumerate}
